\chapter{Repository Structure and Commands}
\label{chap:appendix}

\section{Repository layout}

\begin{itemize}[leftmargin=1.4em]
  \item \texttt{data/processed/recordings/<set>/}, \texttt{data/processed/datasets/<name>/} ---
    one folder per recording set and per dataset variant. \texttt{src/paths.py} defines every
    path used by the project's scripts.
  \item \texttt{scripts/} --- one script per pipeline step, each independently runnable.
  \item \texttt{models/} --- \texttt{registry.json} plus the 12 classical \texttt{.joblib}
    models; \texttt{models/cnn/} holds the CNN weights (model classes live in
    \texttt{scripts/train\_cnn.py}, kept self-contained so it can be copied to a GPU machine).
  \item \texttt{app.py} + \texttt{src/app\_*.py} --- the Dash application.
  \item \texttt{docs/} --- one Markdown report per topic.
\end{itemize}

\section{Key commands}

\begin{verbatim}
python app.py
    # run the application

python scripts/build_dataset.py --annotations official
    # rebuild the official 21-recording dataset

python scripts/train_registry.py
    # train the 12 classical models

python scripts/export_cnn_data.py
    # package data for GPU training -> data/cnn/ (never committed)

python scripts/evaluate_cnn.py
    # score the CNNs (reads models/cnn/ and outputs/predictions/)

python scripts/ensemble_predictions.py --models <key> <key>
    # blend already-trained models, no retraining
\end{verbatim}

\section{Dependencies}

Pinned versions, from \texttt{requirements.txt}:

\begin{multicols}{3}
\small
\begin{itemize}[leftmargin=1.2em,itemsep=0pt]
  \item numpy 2.5.3
  \item scipy 1.18.1
  \item pandas 3.0.6
  \item pyarrow 25.0.1
  \item scikit-learn 1.9.1
  \item joblib 1.6.0
  \item mat-io 1.0.0
  \item h5py 3.16.0
  \item openpyxl 3.1.5
  \item XlsxWriter 3.2.9
  \item PyYAML 6.0.3
  \item matplotlib 3.11.2
  \item seaborn 0.13.2
  \item timescoring 0.0.7
  \item dash 4.4.1
  \item dash-mantine-components 2.8.0
  \item dash-iconify 0.1.2
  \item plotly 6.9.0
  \item plotly-resampler 0.11.1
  \item selenium 4.50.0 (optional)
\end{itemize}
\end{multicols}

No PyTorch is listed: CNN training takes place on separate GPU hardware, and only the resulting
weights and predictions are brought back into this repository.
